\documentclass[11pt]{article}

\usepackage[T1]{fontenc}
\usepackage[utf8]{inputenc}
\usepackage{lmodern}
\usepackage{amsmath,amssymb,amsthm,mathtools}
\usepackage[margin=1.15in]{geometry}
\usepackage{microtype}
\usepackage{enumitem}
\usepackage{xcolor}
\usepackage[colorlinks=true,linkcolor=blue!50!black,citecolor=blue!50!black,urlcolor=blue!50!black]{hyperref}

\newcommand{\R}{\mathbb{R}}
\newcommand{\N}{\mathbb{N}}
\newcommand{\eps}{\varepsilon}
\newcommand{\one}{\mathbf 1}
\newcommand{\abs}[1]{\left\lvert #1\right\rvert}
\newcommand{\set}[1]{\left\{#1\right\}}
\newcommand{\mc}[1]{\mathcal{#1}}
\newcommand{\dd}{\,d}

\theoremstyle{plain}
\newtheorem{theorem}{Theorem}[section]
\newtheorem{proposition}[theorem]{Proposition}
\newtheorem{corollary}[theorem]{Corollary}
\newtheorem{lemma}[theorem]{Lemma}

\theoremstyle{definition}
\newtheorem{definition}[theorem]{Definition}

\theoremstyle{remark}
\newtheorem{remark}[theorem]{Remark}

\title{A Weighting Argument for Pointwise Convergence\\to Convergence in Measure}
\author{Henry Shin}
\date{\today}

\begin{document}

\maketitle

\begin{abstract}
Let $(f_n)$ be a sequence of real-valued Lebesgue measurable functions on
$\R$ such that $f_n(x)\to 0$ for every $x\in\R$.  We give a detailed proof
that there exists a strictly positive measurable function $h$ such that
$f_nh\to 0$ in measure.  The main proof is an expanded version of a
Math StackExchange answer of the author: it uses a finite-measure exhaustion,
Egorov's theorem on each finite-measure piece, and a pointwise-boundedness
stratification.  The construction actually proves the stronger conclusion
that $f_nh\to 0$ uniformly on a conull measurable subset of $\R$.
We also include a shorter dominated-convergence proof of the original
statement and record the corresponding sigma-finite version.
\end{abstract}

\section{Introduction}

On a finite measure space, almost everywhere convergence implies convergence
in measure.  On an infinite measure space, such as $(\R,m)$ with Lebesgue
measure, this implication is false in general.  The problem considered here
asks whether one can multiply by a positive measurable weight in order to
force convergence in measure.

The answer is yes.  More precisely, if $f_n:\R\to\R$ are measurable and
$f_n(x)\to 0$ pointwise, then there is a strictly positive measurable
function $h:\R\to(0,\infty)$ such that
\[
        f_nh\longrightarrow 0
\]
in measure.  The construction below proves something stronger: $h$ can be
chosen so that $f_nh\to 0$ uniformly on a measurable set whose complement has
Lebesgue measure zero.

The key idea is to combine two elementary consequences of pointwise
convergence.  First, after decomposing $\R$ into countably many finite-measure
sets, Egorov's theorem gives countably many large pieces on which the
convergence is uniform.  Second, at each fixed point $x$, the scalar sequence
$(f_n(x))_n$ is bounded, so $\R$ is the union of the level sets on which the
pointwise envelope $\sup_n |f_n|$ is bounded.  The weight $h$ is then made
small on the intersections of these two decompositions.

Throughout this note, convergence in measure on $\R$ means the global
condition
\[
        m\bigl(\{x\in\R: |g_n(x)|>\eps\}\bigr)\longrightarrow 0
        \qquad\text{for every }\eps>0.
\]
Equivalently, one may use $\geq \eps$ instead of $>\eps$.

\section{The strengthened result on \texorpdfstring{$\R$}{R}}

We first state the stronger form proved by the weighting construction.

\begin{theorem}[Conull uniformization by a positive weight]\label{thm:main-R}
Let $f_n:\R\to\R$ be Lebesgue measurable functions such that
\[
        f_n(x)\to 0
        \qquad\text{for every }x\in\R.
\]
Then there exist a measurable set $A\subseteq\R$ and a strictly positive
finite-valued measurable function $h:\R\to(0,\infty)$ such that
\[
        m(\R\setminus A)=0
\]
and
\[
        \sup_{x\in A}|f_n(x)h(x)|\longrightarrow 0.
\]
In particular, $f_nh\to 0$ in measure on $\R$.
\end{theorem}

The proof is split into two organizing lemmas.

\begin{lemma}[Countably many uniform-convergence pieces]\label{lem:egorov-pieces}
There exist measurable sets $A_1,A_2,\dots\subseteq\R$ such that
\[
        m\left(\R\setminus \bigcup_{i=1}^\infty A_i\right)=0,
\]
and such that $f_n\to 0$ uniformly on each $A_i$.
\end{lemma}

\begin{proof}
Choose a finite-measure exhaustion of $\R$; for example, take
\[
        E_r=[-r,r],\qquad r\geq 1.
\]
For each pair $r,k\geq 1$, Egorov's theorem applied on $E_r$ gives a
measurable set $A_{r,k}\subseteq E_r$ such that
\[
        m(E_r\setminus A_{r,k})<\frac{1}{k}
\]
and $f_n\to 0$ uniformly on $A_{r,k}$.

For fixed $r$, we have
\[
        E_r\setminus \bigcup_{k=1}^\infty A_{r,k}
        =\bigcap_{k=1}^\infty (E_r\setminus A_{r,k}),
\]
and hence, for every $k$,
\[
        m\left(E_r\setminus \bigcup_{k=1}^\infty A_{r,k}\right)
        \leq m(E_r\setminus A_{r,k})<\frac{1}{k}.
\]
Therefore
\[
        m\left(E_r\setminus \bigcup_{k=1}^\infty A_{r,k}\right)=0.
\]
Let
\[
        A:=\bigcup_{r,k\geq 1}A_{r,k}.
\]
Since $\R=\bigcup_r E_r$, it follows that
\[
        \R\setminus A
        \subseteq
        \bigcup_{r=1}^\infty
        \left(E_r\setminus \bigcup_{k=1}^\infty A_{r,k}\right),
\]
so $m(\R\setminus A)=0$.  Finally, reindex the countable family
$\{A_{r,k}:r,k\geq 1\}$ as $A_1,A_2,\dots$.  Each reindexed set is measurable,
and $f_n\to 0$ uniformly on each one.
\end{proof}

\begin{lemma}[Pointwise boundedness stratification]\label{lem:boundedness-pieces}
Define
\[
        M(x):=\sup_{n\geq 1}|f_n(x)|.
\]
Then $M$ is measurable and finite everywhere.  Consequently, the sets
\[
        B_j:=\{x\in\R:M(x)\leq j\},\qquad j\geq 1,
\]
are measurable and satisfy
\[
        \R=\bigcup_{j=1}^\infty B_j.
\]
Moreover, on $B_j$ one has
\[
        |f_n(x)|\leq j
        \qquad\text{for every }n\geq 1.
\]
\end{lemma}

\begin{proof}
Since each $f_n$ is measurable, so is $|f_n|$, and therefore the countable
supremum
\[
        M=\sup_{n\geq 1}|f_n|
\]
is measurable.  For each fixed $x\in\R$, the numerical sequence
$(f_n(x))_{n\geq 1}$ converges to $0$, hence is bounded.  Thus $M(x)<\infty$
for every $x$.  The remaining assertions follow immediately from the
definition of $B_j$.
\end{proof}

\begin{proof}[Proof of Theorem \ref{thm:main-R}]
Let $A=\bigcup_i A_i$, where the sets $A_i$ come from
Lemma \ref{lem:egorov-pieces}.  Let $B_j$ be the measurable sets from
Lemma \ref{lem:boundedness-pieces}.  Define
\[
        h(x)
        :=
        \sum_{i,j\geq 1}\frac{1}{j2^{ij}}\one_{A_i\cap B_j}(x)
        +\one_{\R\setminus A}(x).
\]
This is a nonnegative measurable function, being a countable sum of
nonnegative measurable functions.

We first check that $h$ is finite-valued and strictly positive.  Since
\[
        \sum_{i,j\geq 1}\frac{1}{j2^{ij}}
        \leq
        \sum_{i,j\geq 1}\frac{1}{2^{ij}}
        =
        \sum_{i=1}^\infty \frac{1}{2^i-1}
        \leq
        \sum_{i=1}^\infty \frac{1}{2^{i-1}}
        =2,
\]
the defining series is bounded above by $2$ everywhere.  Hence $h$ is finite.
If $x\in\R\setminus A$, then $h(x)\geq 1$.  If $x\in A$, then
$x\in A_i$ for some $i$, and since $\R=\bigcup_j B_j$, also $x\in B_j$ for
some $j$.  The corresponding summand is then positive.  Thus
$h(x)>0$ for every $x\in\R$.

It remains to prove uniform convergence on $A$.  Fix $\eps>0$.  Because
\[
        \sum_{i,j\geq 1}2^{-ij}<\infty,
\]
we may choose $s\geq 1$ so large that
\[
        \sum_{\max\{i,j\}>s}2^{-ij}<\frac{\eps}{2}.
\]
For $x\in A$, the term $\one_{\R\setminus A}(x)$ vanishes, and therefore
\[
\begin{aligned}
        |f_n(x)h(x)|
        &\leq
        \sum_{i,j\geq 1}
        \frac{|f_n(x)|}{j2^{ij}}\one_{A_i\cap B_j}(x) \\
        &=
        \sum_{1\leq i,j\leq s}
        \frac{|f_n(x)|}{j2^{ij}}\one_{A_i\cap B_j}(x)
        +
        \sum_{\max\{i,j\}>s}
        \frac{|f_n(x)|}{j2^{ij}}\one_{A_i\cap B_j}(x).
\end{aligned}
\]
On $A_i\cap B_j$, Lemma \ref{lem:boundedness-pieces} gives
$|f_n(x)|\leq j$.  Hence the tail is bounded uniformly in $x\in A$ and
$n$ by
\[
        \sum_{\max\{i,j\}>s}2^{-ij}<\frac{\eps}{2}.
\]
For the finite block, we have
\[
\begin{aligned}
        \sup_{x\in A}
        \sum_{1\leq i,j\leq s}
        \frac{|f_n(x)|}{j2^{ij}}\one_{A_i\cap B_j}(x)
        &\leq
        \sum_{1\leq i,j\leq s}
        \frac{1}{j2^{ij}}\sup_{x\in A_i}|f_n(x)|.
\end{aligned}
\]
For each fixed $i$, $f_n\to 0$ uniformly on $A_i$, and the sum above has only
finitely many indices.  Therefore the finite block tends to $0$ uniformly on
$A$.  Thus there exists $N$ such that for all $n\geq N$,
\[
        \sup_{x\in A}
        \sum_{1\leq i,j\leq s}
        \frac{|f_n(x)|}{j2^{ij}}\one_{A_i\cap B_j}(x)
        <\frac{\eps}{2}.
\]
Combining the finite-block estimate with the tail estimate gives
\[
        |f_n(x)h(x)|<\eps
        \qquad\text{for all }x\in A\text{ and all }n\geq N.
\]
Hence
\[
        \sup_{x\in A}|f_n(x)h(x)|\to 0.
\]

Finally, since $m(\R\setminus A)=0$, uniform convergence on $A$ implies
convergence in measure.  Indeed, given $\eps>0$, for all sufficiently large
$n$,
\[
        \{x\in\R:|f_n(x)h(x)|>\eps\}
        \subseteq \R\setminus A,
\]
and the right-hand side has measure zero.  Therefore $f_nh\to 0$ in measure.
\end{proof}

\section{A bookkeeping-light version of the same construction}

The double series in the preceding proof is one natural way to combine the
Egorov pieces $A_i$ and the boundedness pieces $B_j$.  A slightly cleaner
presentation is obtained by first making the intersections disjoint.

Enumerate all pairs $(i,j)\in\N^2$ as
\[
        (i_1,j_1),(i_2,j_2),(i_3,j_3),\dots,
\]
and set
\[
        C_p:=A_{i_p}\cap B_{j_p}.
\]
Now disjointify these sets by defining
\[
        D_1=C_1,
        \qquad
        D_p=C_p\setminus \bigcup_{q<p}C_q
        \quad(p\geq 2).
\]
Then the $D_p$ are measurable, pairwise disjoint, and
\[
        A=\bigcup_{p=1}^\infty D_p.
\]
Moreover, if $x\in D_p$, then $x\in A_{i_p}$ and $x\in B_{j_p}$, so
$f_n\to 0$ uniformly on $D_p$ and
\[
        |f_n(x)|\leq j_p
        \qquad (x\in D_p,\ n\geq 1).
\]
Define instead
\[
        \widetilde h(x)
        :=
        \sum_{p=1}^\infty \frac{1}{p j_p}\one_{D_p}(x)
        +\one_{\R\setminus A}(x).
\]
Because the $D_p$ are disjoint, this simply says that
\[
        \widetilde h(x)=\frac{1}{p j_p}\quad\text{if }x\in D_p,
        \qquad
        \widetilde h(x)=1\quad\text{if }x\notin A.
\]
Thus $\widetilde h$ is measurable, finite-valued, and strictly positive.

The proof of uniform convergence is now particularly transparent.  Fix
$\eps>0$ and choose $P$ so large that $1/P<\eps$.  If $p>P$ and $x\in D_p$,
then for every $n$,
\[
        |f_n(x)\widetilde h(x)|
        \leq j_p\cdot \frac{1}{p j_p}
        =\frac{1}{p}
        <\eps.
\]
On the finite union $D_1\cup\cdots\cup D_P$, uniform convergence follows from
uniform convergence on the corresponding finitely many sets $A_{i_p}$.  More
explicitly, for each $p\leq P$, choose $N_p$ so that
\[
        \sup_{x\in D_p}|f_n(x)|<\eps p j_p
        \qquad(n\geq N_p).
\]
With $N=\max_{1\leq p\leq P}N_p$, we get for every $n\geq N$ and every
$x\in D_p$, $p\leq P$,
\[
        |f_n(x)\widetilde h(x)|
        <\eps.
\]
Together with the tail estimate for $p>P$, this proves
\[
        \sup_{x\in A}|f_n(x)\widetilde h(x)|\leq \eps
\]
for all sufficiently large $n$.  Hence $f_n\widetilde h\to 0$ uniformly on
$A$, and therefore in measure.

\begin{remark}
This disjointified proof is the same idea as the double-series proof.  The
sets $A_i$ encode where convergence is already uniform; the sets $B_j$ encode
where the whole sequence is uniformly bounded; and the coefficients of the
weight shrink sufficiently fast as one moves through the countable list of
pieces.
\end{remark}

\section{A shorter proof of the original convergence-in-measure statement}

The stronger conull-uniform result is useful conceptually, but the original
question also has a concise dominated-convergence proof.

Let
\[
        M(x):=\sup_{n\geq 1}|f_n(x)|.
\]
As above, $M$ is measurable and finite everywhere.  Set
\[
        g(x):=e^{-|x|}
\]
and define
\[
        h_0(x):=\frac{g(x)}{1+M(x)}.
\]
Then $h_0$ is measurable and strictly positive.  Moreover,
\[
        |f_n(x)h_0(x)|
        =|f_n(x)|\frac{g(x)}{1+M(x)}
        \leq g(x).
\]
Since $g\in L^1(\R)$ and $f_n(x)h_0(x)\to 0$ pointwise, the dominated
convergence theorem gives
\[
        \int_{\R}|f_nh_0|\,dm\longrightarrow 0.
\]
Thus $f_nh_0\to 0$ in $L^1(\R)$, and hence in measure, since Markov's
inequality gives
\[
        m\bigl(\{x\in\R:|f_n(x)h_0(x)|>\eps\}\bigr)
        \leq
        \frac{1}{\eps}\int_{\R}|f_nh_0|\,dm
        \longrightarrow 0
\]
for every $\eps>0$.

\begin{remark}
The dominated-convergence proof is shorter, but it proves a different kind of
strengthening: $L^1$ convergence rather than uniform convergence on a conull
set.  The Egorov-based construction is more structural.  It shows that after
multiplication by a positive measurable weight, the convergence can be made
uniform away from a null set, even though the underlying measure space is
infinite.
\end{remark}

\section{The sigma-finite form}

Nothing in the Egorov argument is special to the topology of $\R$.  The
relevant measure-theoretic hypothesis is sigma-finiteness.

\begin{theorem}[Sigma-finite version]\label{thm:sigma-finite}
Let $(X,\mc A,\mu)$ be a sigma-finite measure space, and let
$f_n:X\to\R$ be measurable functions such that $f_n(x)\to 0$ for
$\mu$-almost every $x\in X$.  Then there exist a measurable set $A\subseteq X$
with $\mu(X\setminus A)=0$ and a strictly positive finite-valued measurable
function $h:X\to(0,\infty)$ such that
\[
        \sup_{x\in A}|f_n(x)h(x)|\longrightarrow 0.
\]
Consequently $f_nh\to 0$ in measure.
\end{theorem}

\begin{proof}
Let
\[
        X_0:=\{x\in X:f_n(x)\to 0\}.
\]
This set is measurable, since
\[
        X_0=
        \bigcap_{\ell=1}^\infty
        \bigcup_{N=1}^\infty
        \bigcap_{n\geq N}\{x\in X:|f_n(x)|<1/\ell\},
\]
and it is conull by hypothesis.  Choose a countable cover
$X=\bigcup_{r=1}^\infty E_r$ with $\mu(E_r)<\infty$.  Applying Egorov's
theorem on each finite-measure set $E_r\cap X_0$, exactly as in
Lemma \ref{lem:egorov-pieces}, gives countably many measurable sets
$A_1,A_2,\dots\subseteq X_0$ such that
\[
        \mu\left(X\setminus \bigcup_{i=1}^\infty A_i\right)=0,
\]
and $f_n\to 0$ uniformly on each $A_i$.  Put
\[
        A:=\bigcup_{i=1}^\infty A_i.
\]

On $X_0$, define the envelope
\[
        M(x):=\sup_{n\geq 1}|f_n(x)|.
\]
It is finite on $X_0$.  For $j\geq 1$, define
\[
        B_j:=X_0\cap\{x\in X:M(x)\leq j\}.
\]
Then the $B_j$ are measurable, $X_0=\bigcup_j B_j$, and
$|f_n|\leq j$ on $B_j$ for all $n$.

Now define
\[
        h(x)
        :=
        \sum_{i,j\geq 1}\frac{1}{j2^{ij}}\one_{A_i\cap B_j}(x)
        +\one_{X\setminus A}(x).
\]
The same estimates as in the proof of Theorem \ref{thm:main-R} show that
$h$ is measurable, finite-valued, strictly positive everywhere, and that
$f_nh\to 0$ uniformly on $A$.  Since $\mu(X\setminus A)=0$, convergence in
measure follows.
\end{proof}

The shorter $L^1$ proof also has a sigma-finite version.

\begin{proposition}[Sigma-finite dominated-convergence proof]\label{prop:l1-sigma-finite}
Let $(X,\mc A,\mu)$ be sigma-finite, and suppose $f_n:X\to\R$ are measurable
and $f_n\to 0$ almost everywhere.  Then there is a strictly positive
finite-valued measurable function $h$ such that
\[
        \int_X |f_nh|\,d\mu\longrightarrow 0.
\]
In particular, $f_nh\to 0$ in measure.
\end{proposition}

\begin{proof}
Let $X_0=\{x:f_n(x)\to 0\}$, and let
\[
        M(x):=\sup_{n\geq 1}|f_n(x)|.
\]
The function $M$ is an extended-valued measurable function and is finite on
$X_0$.  Define the finite-valued measurable function
\[
        \widetilde M(x):=
        \begin{cases}
        M(x), & x\in X_0,\\
        0, & x\notin X_0.
        \end{cases}
\]

Since $\mu$ is sigma-finite, choose measurable sets $E_r$ with
$X=\bigcup_r E_r$ and $\mu(E_r)<\infty$.  Passing to the disjoint refinement
\[
        F_1=E_1,
        \qquad
        F_r=E_r\setminus\bigcup_{q<r}E_q\quad(r\geq 2),
\]
we obtain a measurable partition of $X$ with $\mu(F_r)<\infty$.  Define
\[
        g(x):=\sum_{r=1}^\infty
        \frac{2^{-r}}{1+\mu(F_r)}\one_{F_r}(x).
\]
Then $g$ is strictly positive, finite-valued, measurable, and integrable,
since
\[
        \int_X g\,d\mu
        =
        \sum_{r=1}^\infty
        \frac{2^{-r}\mu(F_r)}{1+\mu(F_r)}
        \leq
        \sum_{r=1}^\infty 2^{-r}<\infty.
\]
Define
\[
        h(x):=\frac{g(x)}{1+\widetilde M(x)}.
\]
Then $h$ is strictly positive, finite-valued, and measurable.  On the conull
set $X_0$,
\[
        |f_nh|
        \leq g
        \qquad\text{and}\qquad
        f_nh\to 0.
\]
The dominated convergence theorem therefore gives
\[
        \int_X |f_nh|\,d\mu\to 0.
\]
The convergence in measure follows from Markov's inequality.
\end{proof}

\section{Concluding remarks}

The original construction is a useful pattern.  When pointwise convergence is
too weak globally, one can often separate the problem into two countable
stratifications:
\begin{enumerate}[label=\textup{(\roman*)},leftmargin=2.5em]
    \item pieces on which convergence is uniform, supplied here by Egorov's
    theorem on finite-measure pieces;
    \item pieces on which the whole sequence is uniformly bounded, supplied
    here by the pointwise envelope $M=\sup_n |f_n|$.
\end{enumerate}
The weight is then chosen small enough on the intersections of the two
stratifications.  The role of sigma-finiteness is precisely to make the first
stratification countable through finite-measure exhaustion.  The role of
pointwise convergence is not only to give the limiting value $0$, but also to
guarantee pointwise boundedness of the sequence, which is what makes the
second stratification possible.

\begin{thebibliography}{9}

\bibitem{MSE}
Henry Shin,
\emph{Answer to: Let $f_n$ be real and measurable with $f_n(x)\to 0$; show
there is a positive measurable function $h$ such that $f_nh\to 0$ in measure},
Mathematics Stack Exchange, answered July 1, 2021.
\url{https://math.stackexchange.com/questions/4186426/let-f-n-be-real-measurable-w-f-nx-rightarrow-0-show-there-is-a-positiv}

\bibitem{folland}
Gerald B. Folland,
\emph{Real Analysis: Modern Techniques and Their Applications},
2nd ed., Wiley, 1999.

\bibitem{royden}
H. L. Royden and P. M. Fitzpatrick,
\emph{Real Analysis},
4th ed., Pearson, 2010.

\end{thebibliography}

\end{document}
